% Mechanism illustration: all data generated by the restored original C++ validator.
% Source figures: mechanisms/figures/*.pdf. Main manuscript insertion remains with the parent editor.

\begin{figure*}[t]
\centering
\includegraphics[width=.97\textwidth]{figures/observed_continuation_mechanism.pdf}
\caption{Observed-space dependence of terminal continuation. (a,b) Two validation maps at $t=0.8$ s after five lidar scans, with identical aircraft state, goal, primary P69 and 15-control successor bank. The background shows the two-metre voxel slice $z\in[20,22)$ m; curves are horizontal projections of three-dimensional rollouts. P69 is primary-feasible in both maps, but admits no valid successor in Map A and one in Map B. P107 is an alternative primary with two valid successors in Map A; its first valid successor is shown. Crosses mark rejected rollout endpoints. (c) Full-bank outcomes from the three-dimensional validator, including occupancy, unknown-space and flight-envelope checks. (d) Terminal free-space probe fractions for P69, grouped by azimuth with elevations $-20^{\circ},0^{\circ},20^{\circ}$.}
\label{fig:observed_mechanism}
\end{figure*}

% Suggested method/experiment paragraph, approximately 100 words.
Figure~\ref{fig:observed_mechanism} separates primary feasibility from terminal continuation. In Map A, both P69 and P107 pass primary validation, while only P107 admits a successor in the finite bank. Applying the same P69 rollout and successor controls in Map B changes its continuation label from zero to one. The aircraft state, goal and all 28 state--task features are identical across the two inputs; 18 of the 30 observation features differ. Thus, a state-only predictor cannot distinguish these outcomes. The observation probes encode the local free-space variation that changes continuation feasibility.

% Optional one-sentence quantitative support if room permits.
The validation set contains 1,186 state--task feature groups with opposite continuation labels; every group includes a pair with identical recorded aircraft state and primary-library index.

\begin{figure*}[t]
\centering
\includegraphics[width=.97\textwidth]{figures/urban_query_trace.pdf}
\caption{Candidate scheduling and incumbent-cost pruning on the first prescribed urban input whose seed candidate has no continuation witness ($t=2$ s). (a) Each number is the candidate's one-based rank in stable task-cost order. Budgeted schedules use $K=12$; hatched entries are pruned before successor integration and receive no replacement. Cost and uncapped Full stop at the witness of rank 7. (b) Cumulative successor integration steps, including the common seed query. Selected and Diverse have identical schedules on this input. Obs+Aux verifies four banks and prunes eight candidates; Selected verifies six banks and prunes six.}
\label{fig:query_trace}
\end{figure*}

% Suggested algorithm/experiment paragraph, approximately 100 words.
Figure~\ref{fig:query_trace} shows how scheduling and pruning allocate work on one urban input. Selected and Diverse find a witness at the sixth scheduled candidate and prune all six remaining entries, using 2,413 integration steps. Obs+Aux first finds a witness at slot three and improves it at slot five; the intervening and remaining higher-cost candidates are skipped. It verifies four banks but uses 3,080 steps because its rejected rollouts are longer. Cost and Full reach the rank-7 witness after 7,382 steps. The trace illustrates why candidate counts and measured integration work capture different aspects of the validation budget.
